\documentclass[11pt, a4paper]{article}

\usepackage[T1]{fontenc}
\usepackage[utf8]{inputenc}
\usepackage{tgpagella}
\usepackage{tgcursor}
\usepackage{microtype}
\usepackage[spanish, es-noshorthands]{babel}
\usepackage{amsmath, amssymb}
\usepackage[table, xcdraw]{xcolor}
\usepackage{booktabs}
\usepackage{tabularx}
\usepackage[margin=1.5cm]{geometry}
\usepackage{fancyhdr}

\pagestyle{fancy}
\fancyhf{}
\setlength{\headheight}{16pt}
\lhead{\textbf{IMT-342} | Worksheet Validación URDF}
\rhead{Hito 1 - Semana 6[cite: 8]}

\begin{document}

\section*{A. Identificación[cite: 8]}
\textbf{Grupo:} \rule{4cm}{0.15mm} \quad \textbf{URDF File/Versión:} \rule{4cm}{0.15mm} \\
\textbf{Integrantes:} \rule{10cm}{0.15mm} \\
\textbf{Base Frame:} \rule{4cm}{0.15mm} \quad \textbf{Terminal/Flange Frame:} \rule{4cm}{0.15mm}

\section*{B. Validez XML / URDF[cite: 8]}
\textbf{Herramienta/Método:} \rule{4cm}{0.15mm} \quad \textbf{Resultado:} Pass / Fail \\
\textbf{Error y Corrección:} \rule{10cm}{0.15mm}

\section*{C. Kinematic Tree[cite: 8]}
\textit{(Dibuje el árbol jerárquico)}
\vspace{1.5cm}

Checklist: $\square$ Conexiones completas $\square$ Cero ramas accidentales $\square$ Nombres exactos $\square$ Flange identificable

\section*{D. Evidencia de Definición de Juntas[cite: 8]}
\begin{center}
    \footnotesize
    \renewcommand{\arraystretch}{1.5}
    \noindent\begin{tabularx}{\textwidth}{|c|c|c|X|X|c|c|X|}
    \hline
    \textbf{Joint} & \textbf{Parent} & \textbf{Child} & \textbf{\texttt{xyz}} & \textbf{\texttt{rpy}} & \textbf{\texttt{axis}} & \textbf{Límite Fuente} & \textbf{Justificación Física} \\ \hline
    J1 & & & & & & & \\ \hline
    J2 & & & & & & & \\ \hline
    J3 & & & & & & & \\ \hline
    J4 & & & & & & & \\ \hline
    J5 & & & & & & & \\ \hline
    J6 & & & & & & & \\ \hline
    \end{tabularx}
\end{center}

\section*{E. Zero / Reference Configuration[cite: 8]}
$q^{(0)} = [ \quad, \quad, \quad, \quad, \quad, \quad ]^T$ \\
\textbf{Pose esperada cualitativamente:} \rule{8cm}{0.15mm} \\
\textbf{¿Pose cero suficiente para probar corrección? ¿Por qué?} \rule{6cm}{0.15mm}

\section*{F. Pruebas de Juntas Individuales[cite: 8]}
\begin{center}
    \footnotesize
    \noindent\begin{tabularx}{\textwidth}{|c|X|c|X|X|c|c|}
    \hline
    \textbf{Joint} & \textbf{Demás fijas en} & \textbf{Test Val} & \textbf{Mov. Esperado} & \textbf{Mov. Observado} & \textbf{Match?} & \textbf{Notas} \\ \hline
    J1 & & & & & $\square$ & \\ \hline
    J2 & & & & & $\square$ & \\ \hline
    J3 & & & & & $\square$ & \\ \hline
    J4 & & & & & $\square$ & \\ \hline
    J5 & & & & & $\square$ & \\ \hline
    J6 & & & & & $\square$ & \\ \hline
    \end{tabularx}
\end{center}

\section*{G \& H. Configuraciones de Prueba 1 y 2 (Comparación FK)[cite: 8]}
\textbf{Test A:} $q^{(A)} = [ \quad, \quad, \quad, \quad, \quad, \quad ]^T$. Ref. Matemática $p^{(A)} =$ \rule{2cm}{0.15mm}, $R^{(A)} =$ \rule{2cm}{0.15mm}. \\
Evidencia ROS/TF: \rule{10cm}{0.15mm} \\
\textbf{Test B:} $q^{(B)} = [ \quad, \quad, \quad, \quad, \quad, \quad ]^T$. Ref. Matemática $p^{(B)} =$ \rule{2cm}{0.15mm}, $R^{(B)} =$ \rule{2cm}{0.15mm}. \\
Comparación: $\square$ Coherente \quad $\square$ Discrepancia \quad $\square$ Pendiente relación de marcos.

\section*{I. Discrepancia y Debugging[cite: 8]}
\textbf{Síntoma:} \rule{12cm}{0.15mm} \\
\textbf{Revisado:} $\square$ Sintaxis $\square$ Parent/Child $\square$ Origin xyz/rpy $\square$ Axis $\square$ Visual origin $\square$ ROS $\square$ DH vs URDF \\
\textbf{Corrección y evidencia:} \rule{12cm}{0.15mm}

\section*{J. Conclusión Final y Preparación Individual[cite: 8]}
\textbf{Grupo:} \textit{Consideramos que nuestro URDF es cinemáticamente defendible porque...} \rule{5cm}{0.15mm} \\
\textbf{Individual:} \textit{(Cada miembro preparará su defensa individual respondiendo: ¿Qué junta defenderé? ¿Por qué su axis/origin son correctos? ¿Qué diferencia existe entre DH y URDF en esta junta?)}

\end{document}
